The target is a network of known total length, not the whole two-dimensional
hexagon. Each triangle has a local upper bound on the length it can cover.
Because the covering sets are open, some covering pieces must overlap; hence
the sum of their lengths must be \emph{strictly greater} than the target
length. An upper bound equal to that target is already a contradiction.

For the skeleton, an ordinary nonsupercritical V triangle without adjacent
radial support contributes at most two. Supercritical V triangles and those
with adjacent radial support each contribute less than $3/2$. When the
C triangle reaches the boundary, its skeleton contribution is also below
$3/2$. The count argument below simply adds these local budgets. The local
caps require their stated hypotheses; no unrestricted skeleton cap is being
asserted for the C triangle.
